\section{The Auditor: Stage One of Three}
\label{sec:auditor}

Figure~\ref{fig:pipeline} states the pipeline and, for each stage, what it does and does not
license; conflating those stages is what produces overclaims. This section is stage 1.

Tiers 1--3 are properties of the corpus, so they can be screened with no model, no GPU and no
forward pass. Given a corpus of equivalence classes the auditor reports, per tier, whether class
representatives are separable by that tier's features alone: the \emph{unique var-set fraction}
(tier 1: if high, knowing which variables appear pins the class before any model exists),
\emph{operator-bag separability} with its collision profile (tier 2), the \emph{arity distribution}
(which often partitions a small suite into near-singleton groups, a confound aggregate accuracy
hides), and the tier-3 statistics of \S\ref{sec:tier3}. Figure~\ref{fig:cli} (Appendix~\ref{app:cli}) shows it on two corpora;
the tool is dependency-light and takes a documented JSON format, with adapters for every corpus in
Table~\ref{tab:benchmark_survey}.

\paragraph{What it does not do.} These columns measure whether class \emph{representatives} differ;
they do not predict what a baseline scores on \emph{held-out forms}. Across 12 corpora no
zero-parameter separability statistic we tried predicts untrained-encoder accuracy ($R^2 = -2.8$;
Appendix~\ref{app:screen}), reported as a limitation rather than tuned until it correlated.
Separability is a \emph{diagnostic}; measured baseline performance is the \emph{evidentiary}
criterion. A cheap pre-training screen that tells you where to look is worth having even when it
cannot tell you how much the cue will be worth.
